% Copyright (c) 2026 Eugene Joseph M. Ragasa.
% SPDX-License-Identifier: Apache-2.0
% Author: Eugene Joseph M. Ragasa
% Formal proof completed: 2026-10-19
\documentclass[11pt]{article}
\usepackage[T1]{fontenc}\usepackage{amsmath,amssymb}\usepackage[margin=1in]{geometry}\usepackage[hidelinks]{hyperref}
\title{Proof Dossier: Canonical Pauli Decomposition}\author{Eugene Joseph M. Ragasa}\date{Formal proof completed: 2026-10-19}
\begin{document}\maketitle
\section{Identity and status}
ID \texttt{BG-FOUNDATION-THM-007}; status \texttt{proved}. Lean declaration:
\par\noindent {\tiny\path{ProjectKoios.BandGauge.Foundations.}}\\
{\tiny\path{eq_rankTwoPauliMatrix_of_conjTranspose_eq_self_of_trace_eq_zero}}.
\section{Statement and proof}
Every traceless Hermitian complex $2$-by-$2$ matrix equals
$x\sigma_x+y\sigma_y+z\sigma_z$, where
$x=\operatorname{Re}M_{01}$, $y=-\operatorname{Im}M_{01}$, and
$z=\operatorname{Re}M_{00}$. Hermiticity makes the diagonal real and identifies the
lower-left entry with the conjugate upper-right entry; trace zero fixes the second
diagonal entry.
\section{Boundary}
This is standard supporting linear algebra, not a novelty claim. Kernel axioms are
\texttt{propext}, \texttt{Classical.choice}, and \texttt{Quot.sound}; no admissions.
AI-assisted and kernel-checked. No independent expert review is recorded.
\end{document}
